\subsection{例：测度空间中的线性表示}

同样设 G 为局部紧群，在局部紧空间 X 上连续左作用，并令 \(\chi\) 为 \(\mathbf{G} \times \mathbf{X}\) 上的连续乘子。\(\boldsymbol{\gamma}_{\chi}\) 作为 \(\mathbf{G}\) 在 \(\mathcal{K}(\mathbf{X})\) 中的线性表示，在 \(\mathcal{M}(\mathbf{X})\) 中有一个逆伴随表示，我们仍记为 \(\boldsymbol{\gamma}_{\chi}\)，它由下式定义（其中 \(\mu \in \mathcal{M}(\mathbf{X})\)、\(f \in \mathcal{K}(\mathbf{X})\)）：

\[
\langle \boldsymbol {\gamma} _ {\chi} (s) \mu , f \rangle = \langle \mu , \boldsymbol {\gamma} _ {\chi} (s ^ {- 1}) f \rangle = \langle \chi (s, \cdot) \cdot \mu , \boldsymbol {\gamma} (s ^ {- 1}) f \rangle = \langle \boldsymbol {\gamma} (s) (\chi (s, \cdot) \cdot \mu), f \rangle ,
\]

从而

\[
\boldsymbol {\gamma} _ {\chi} (s) \mu = \boldsymbol {\gamma} (s) (\chi (s, \cdot) \cdot \mu) = (\boldsymbol {\gamma} (s) \chi (s, \cdot)) \cdot (\boldsymbol {\gamma} (s) \mu).
\]

注意到

\[
\left(\gamma (s) \chi (s, \cdot)\right) (x) = \chi (s, s ^ {- 1} x).
\]

\(\gamma_{\chi}\) 作为 G 在 \(\mathcal{C}(X)\) 中的线性表示，在由 X 上具有紧支集的测度构成的空间 \(\mathcal{C}'(X)\) 中有一个逆伴随表示，我们仍将该表示记为 \(\gamma_{\chi}\)；\(\gamma_{\chi}(s)\) 的自同态在 \(\mathcal{C}'(X)\) 上，是 \(\gamma_{\chi}(s)\) 的自同态在 \(\mathcal{M}(X)\) 上的限制。

命题 6. --- 若在 \(\mathcal{M}(\mathrm{X})\)（分别在 \(\mathcal{C}'(\mathrm{X})\)）上赋予关于 \(\mathcal{H}(\mathrm{X})\)（分别关于 \(\mathcal{C}(\mathrm{X})\)）的紧子集上一致收敛的拓扑，则 \(\gamma_{\chi}\) 作为 \(\mathrm{G}\) 在 \(\mathcal{M}(\mathrm{X})\)（分别在 \(\mathcal{C}'(\mathrm{X})\)）中的线性表示是连续的。

命题 7. --- 假设每个函数 \(\chi(s,\cdot)\) 都有界。则 \(\gamma_{\chi}\) 保持 \(\mathcal{M}^{1}(X)\) 不变；若在 \(\mathcal{M}^{1}(X)\) 上赋予关于 \(\overline{\mathcal{K}(X)}\) 的紧子集上一致收敛的拓扑，则 \(\gamma_{\chi}\) 作为 \(G\) 在 \(\mathcal{M}^{1}(X)\) 中的线性表示是连续的。

这些命题由命题 3、4、5 得到。

